\documentclass[conference]{IEEEtran}
\IEEEoverridecommandlockouts
\usepackage[T1]{fontenc}
\usepackage[utf8]{inputenc}
\usepackage{amssymb}
\usepackage{xcolor}
\usepackage{listings}
\usepackage{textcomp}
\usepackage{hyperref}
\lstset{
  basicstyle=\ttfamily\footnotesize,
  breaklines=true,
  frame=single,
  columns=fullflexible,
  keepspaces=true,
  showstringspaces=false,
  literate={"}{\textquotedbl}1
}
\begin{document}

\title{Hands-On Computer Vision and Prosthetic Control:\\ A Beginner Workshop}
\author{\IEEEauthorblockN{Severus Prosthetics Workshop}
\IEEEauthorblockA{Motion Recognition Workshop Series}}
\maketitle

\begin{abstract}
This workshop teaches two linked skills in roughly one hour: real-time computer-vision hand tracking, and translating that tracking into control of a prosthetic hand. Participants extend a working program by completing five short edits, each of which switches on one stage of the pipeline -- detection, visualisation, per-finger state estimation, and actuator command mapping. The program runs at every stage, degrading gracefully so unfinished features are simply labelled on screen. This beginner edition requires no prior programming experience: each edit is an uncomment-one-line choice with the answer explained in full.
\end{abstract}

\begin{IEEEkeywords}
computer vision, MediaPipe, hand tracking, prosthetics, human-machine interface, gesture recognition, education
\end{IEEEkeywords}

\section{Introduction}
Modern prosthetic hands increasingly take their commands from sensors that estimate user intent. A camera plus a hand-tracking model is one of the most accessible such sensors: it needs no electrodes and runs on a laptop. This workshop builds that path end to end, from pixels to finger motion, using the Severus hand and its simulator.

The exercise file opens a live window immediately. Five marked lines (named \lstinline!ANSWER_1! through \lstinline!ANSWER_5!) each enable one stage. As you complete them, the program progresses from doing very little, to drawing the hand, to recognising finger states, to moving a simulated (and optionally physical) prosthesis.

\section{Background}
MediaPipe estimates 21 landmarks per hand. Each landmark has normalised image coordinates, where x increases to the right and y increases DOWNWARD from the top of the frame. A finger is treated as extended when its fingertip is higher than the knuckle below it -- a smaller y. The thumb is handled on the x axis instead, because it moves sideways.

These per-finger booleans become commands. In simulation each finger has an openness between 0.0 (closed) and 1.0 (open); the identical command stream drives the physical Severus hand once it is connected. The workshop therefore mirrors a real human-machine interface: sense, interpret, actuate.

\section{Workshop Exercises}

\subsection{Exercise 1: Hand-landmark detection}

\textbf{Line to edit:} \lstinline!ANSWER_1!.

Uncomment the correct option so the line reads:

\begin{lstlisting}
ANSWER_1 = vision.HandLandmarker
\end{lstlisting}


\textbf{What this line does.} Selects the MediaPipe vision task that, for every frame, returns the 21 hand keypoints (knuckles and fingertips) as normalised coordinates. The program then builds its detector from this class.

\textbf{The choices.} Alternatives are deliberately offered: PoseLandmarker tracks the whole body and GestureRecognizer only emits named gestures such as thumbs-up. Neither exposes the per-finger geometry we need.

\textbf{Why it matters for control.} This is the SENSING front end of the prosthesis: it converts a raw camera image into structured numbers. Without a reliable hand estimate there is no signal to map onto the fingers, so every later step depends on this one.

\subsection{Exercise 2: Drawing the landmarks}

\textbf{Line to edit:} \lstinline!ANSWER_2!.

Uncomment the correct option so the line reads:

\begin{lstlisting}
ANSWER_2 = draw_landmarks
\end{lstlisting}


\textbf{What this line does.} Picks the drawing function that overlays the 21 keypoints and the skeleton connections onto the camera frame so the tracked hand is visible.

\textbf{The choices.} The other option, overlay\_states, writes the finger words as text. Useful, but it does not show the geometry that proves tracking works.

\textbf{Why it matters for control.} Visual feedback is how an engineer validates a sensor before trusting it for control. Seeing the skeleton lock onto the hand confirms the estimate is stable enough to drive an actuator.

\subsection{Exercise 3: Finger extend / retract rule}

\textbf{Line to edit:} \lstinline!ANSWER_3!.

Uncomment the correct option so the line reads:

\begin{lstlisting}
ANSWER_3 = "<"
\end{lstlisting}


\textbf{What this line does.} Defines, for the four fingers, the rule that turns two landmark heights into a binary state. The engine evaluates tip-y (operator) pip-y; you supply the operator. Image coordinates put y = 0 at the TOP and increase downward, so a raised (extended) fingertip has the SMALLER y.

\textbf{The choices.} Choosing ">" simply inverts the logic (extended reads as retracted), which is an instructive bug to see on screen.

\textbf{Why it matters for control.} This is the heart of joint-space mapping: a continuous measurement is thresholded into an open/close command per finger. The same idea scales up to proportional control, but binary is the clearest start.

\subsection{Exercise 4: Thumb extend / retract rule}

\textbf{Line to edit:} \lstinline!ANSWER_4!.

Uncomment the correct option so the line reads:

\begin{lstlisting}
ANSWER_4 = ">"
\end{lstlisting}


\textbf{What this line does.} Defines the thumb rule. Because the thumb abducts SIDEWAYS rather than up and down, it is classified on the x axis: the engine evaluates tip-x (operator) ip-x for a right hand and mirrors it automatically for a left hand.

\textbf{The choices.} x = 0 is the LEFT edge and increases to the right, so for a right hand an abducted thumb tip sits at a larger x than its lower joint.

\textbf{Why it matters for control.} Real hands are not uniform; the thumb has a different kinematic axis. Treating it separately mirrors how prosthetic controllers special-case the thumb, which dominates grasp quality.

\subsection{Exercise 5: Driving the simulated hand}

\textbf{Line to edit:} \lstinline!ANSWER_5!.

Uncomment the correct option so the line reads:

\begin{lstlisting}
ANSWER_5 = 1.0
\end{lstlisting}


\textbf{What this line does.} Connects the detected states to the actuator model. Each simulated finger takes an openness value in the range 0.0 to 1.0; the engine uses ANSWER\_5 for an extended finger and one-minus-ANSWER\_5 for a retracted one, then eases smoothly toward that target.

\textbf{The choices.} Setting it to 0.0 inverts the hand (your open hand would close it), a vivid demonstration of why the command convention must be defined.

\textbf{Why it matters for control.} This is the ACTUATOR command mapping. The very same per-finger command stream drives the on-screen hand and, once an Arduino is connected, the physical Severus prosthesis -- simulation and hardware share one interface.

\section{Conclusion}
Completing the five edits yields a full sense-interpret-actuate loop: a camera tracks the hand, per-finger states are inferred, and a prosthetic hand mirrors the motion. The same structure underlies production myoelectric and vision-based controllers; only the sophistication of each stage grows. If a finished file is connected to a prepared Severus hand, it will move the real fingers exactly as it moves the simulated ones.

\end{document}
